\documentclass[11pt]{article}
\usepackage{amsmath, amssymb, amsthm, enumerate}
\usepackage[margin=1in, top=1.2in, bottom=1.2in]{geometry}
\usepackage{parskip}
\usepackage[colorlinks=true, linkcolor=blue, citecolor=blue, urlcolor=blue]{hyperref}
\usepackage{cleveref}
\theoremstyle{definition}
\newtheorem{definition}{Definition}[section]
\newtheorem{theorem}{Theorem}[section]
\newtheorem{lemma}{Lemma}[section]
\newtheorem{proposition}{Proposition}[section]
\newtheorem{corollary}{Corollary}[section]
\theoremstyle{remark}
\newtheorem{remark}{Remark}[section]
\newtheorem{example}{Example}[section]
\setlength{\parindent}{0pt}
\setlength{\parskip}{6pt}
\begin{document}

\begin{center}
{\Large\textbf{Techiniques, Examples And Counterexamples}}\\
\vspace{0.3cm}
{\normalsize\today}
\end{center}
\vspace{0.5cm}


% --- Page 1 ---
\begin{itemize}
    \item \textbf{Modules $(M, R, +)$}
    \begin{itemize}
        \item Verify the axioms
    \end{itemize}

    \item \textbf{Submodules}
    \begin{itemize}
        \item Verify the axioms (Check subgroup and for any $r \in R$, $m \in M$, $rm \in M$.)
    \end{itemize}

    \item \textbf{Homomorphism}
    \begin{itemize}
        \item Verify the axioms
    \end{itemize}
\end{itemize}

% --- Page 2 ---
% blank page

% --- Page 3 ---
\begin{example}\label{ex:modules}
    \begin{enumerate}
        \item Vector space $\mathbb{C}$ (Ring over module is a field).
        \item $(V, T)$ where $V$ is a vector space and $T$ is a linear transform. Ring is constructed to be $F[x]$ where
        \[
        f(x) \cdot v := f(T)(v).
        \]
        Note: $(V, T) \longmapsto F[x]$ module.
        For $\longmapsto$ we define $T(x) := x \cdot v$.
    \end{enumerate}
\end{example}

% --- Page 4 ---
\begin{itemize}
    \item \textbf{Submodules}

    \begin{definition}\label{def:submodule_z}
        In $\mathbb{Z}$-modules, say $M$, $N$ is a submodule if $(N, +) \leq (CM, +)$.
    \end{definition}

    \begin{definition}\label{def:submodule_r}
        In $R$ is a $R$-module, $N$ is a submodule if $N$ is an ideal.
    \end{definition}

    \begin{definition}\label{def:submodule_fx}
        If $M$ is a $F[x]$-module $\cong (V, T)$, $W$ is a submodule if $T(CW) \subseteq W$, i.e., $W$ is invariant under $T$.
    \end{definition}
\end{itemize}

% --- Page 5 ---
\begin{center}
\textbf{Finitely generated and cyclic modules}
\end{center}

\begin{tabular}{|l|l|l|}
\hline
\textbf{Module} & \textbf{f.g.} & \textbf{cyclic} \\ \hline
$(V, \mathbb{F})$ & $\dim V < \infty$ & $\dim V = \{0, 1\}$ \\ \hline
$G$ is abelian group & f.g as & cyclic as \\
treated like $\mathbb{Z}$-module & group & group \\ \hline
$R$ comm ring and $M \triangleleft R$ & f.g as & principle \\
$M$ is $R$-module & an ideal & ideal \\ \hline
\end{tabular}

\begin{definition}\label{def:finitely_generated_module}
A module $M$ is \textit{finitely generated} if there exists a finite set of elements $\{m_1, m_2, \dots, m_n\}$ in $M$ such that every element of $M$ can be written as a linear combination of these elements.
\end{definition}

\begin{definition}\label{def:cyclic_module}
A module $M$ is \textit{cyclic} if there exists an element $m \in M$ such that every element of $M$ can be written as a multiple of $m$.
\end{definition}

\begin{remark}\label{rem:vector_space}
For a vector space $(V, \mathbb{F})$, being finitely generated means $\dim V < \infty$. A cyclic vector space has dimension either $0$ or $1$.
\end{remark}

\begin{remark}\label{rem:abelian_group}
An abelian group $G$ treated as a $\mathbb{Z}$-module is finitely generated if it is finitely generated as a group. It is cyclic as a $\mathbb{Z}$-module if it is cyclic as a group.
\end{remark}

\begin{remark}\label{rem:ideal_module}
For a commutative ring $R$ and an $R$-module $M$ that is an ideal of $R$, $M$ is finitely generated as an $R$-module if it is finitely generated as an ideal. A principal ideal is a cyclic $R$-module.
\end{remark}

% --- Page 6 ---
\begin{itemize}
    \item \textbf{Module homomorphisms}

    \begin{itemize}
        \item Abelian groups $\leftrightarrow$ $\mathbb{Z}$-modules, then group homomorphism $\leftrightarrow$ module homomorphism.
    \end{itemize}

    \begin{itemize}
        \item For vector fields, linear transformation $\leftrightarrow$ module homomorphism.
    \end{itemize}

    \begin{itemize}
        \item For ring $R$, $\Theta: M \rightarrow N$ where $M = N = R$ and
        \[
        \Theta(rs) = r\Theta(s)
        \]
        is \textbf{not} a ring homomorphism but is a module homomorphism.
    \end{itemize}

    \begin{itemize}
        \item Isomorphism theorems hold for modules.
    \end{itemize}
\end{itemize}

% --- Page 7 ---
% blank page

\end{document}